% Tělo sekce: Generování a úprava obrázků a grafiky
Tato sekce shrnuje praktické možnosti tvorby i editace vizuálů přímo v rámci nástrojů Microsoft (zejména Designer) a ukazuje doporučené pracovní postupy a ukázkové prompty pro použití ve výuce. Uváděné technické možnosti vycházejí z dostupných funkcí Designeru a z funkcí pro editaci obrázků v Copilot chatu \cite{kb_ai_101_tipu_pdf_04e4a115_page_51_1,kb_ai_101_tipu_pdf_04e4a115_page_52_2}.

Krátké shrnutí schopností Designeru (co využijete ve výuce)
\begin{itemize}
  \item Designer pracuje ve třech režimech: generování nové grafiky z promptu, úprava stávající grafiky pomocí AI a klasický ruční návrh \cite{kb_ai_101_tipu_pdf_04e4a115_page_51_1}.
  \item Umí generovat obrázky na základě textového promptu a umožňuje definovat jiný než čtvercový poměr stran (rychle vytvoříte širokoúhlé bannery pro slide) \cite{kb_ai_101_tipu_pdf_04e4a115_page_51_1}.
  \item Podporuje také specifické kategorie užitečné pro školní použití: monogramy, omalovánky, avatary a nálepky (sticker designy včetně bílého okraje) \cite{kb_ai_101_tipu_pdf_04e4a115_page_51_1}.
  \item Editace části existujícího obrázku (např. změna pozadí, úprava nápisu nebo přidání/odebrání doplňků) je možná prostřednictvím funkcí transformace v Copilot chatu — tato funkce vyžaduje placený Microsoft 365 Copilot \cite{kb_ai_101_tipu_pdf_04e4a115_page_52_2}.
\end{itemize}

Doporučený rychlý workflow pro vytvoření vizuálu k lekci
\begin{enumerate}
  \item Definujte účel obrázku: (a) slide ilustrace, (b) cover pro materiál, (c) ikonka/štítek pro aktivitu.
  \item Zvolte formát a poměr stran podle cílového použití (např. širokoúhlý pro slide 16:9, čtverec pro ikonu, vertikál pro story/banner). (praktické doporučení; viz možnosti Designeru \cite{kb_ai_101_tipu_pdf_04e4a115_page_51_1})
  \item Napište krátký prompt s požadovaným stylem, barvami a elementy (viz šablony níže).
  \item Vygenerujte 3 varianty, vyberte nejvhodnější a případně použijte editaci části obrázku (transformace) pro doladění \cite{kb_ai_101_tipu_pdf_04e4a115_page_52_2}.
  \item Exportujte ve vhodném formátu (PNG/SVG pro ikony, JPG/PDF pro materiály) a uložte originální prompt/metadata pro pozdější reprodukci. (general background knowledge)
\end{enumerate}

Praktické prompt‑šablony (copy \\ \& adapt)
\begin{verbatim}
1) Monogram pro kurz:
"Vytvoř monogram pro kurz 'Fyzika 1' v jednoduchém moderním stylu,
barvy: tmavě modrá + světle zelená, čistý vektorový styl, výstup 1:1,
jednoduchý obrys vhodný pro tisk na samolepku."
2) Motivace na začátek hodiny (wide banner 16:9):
"Motivační banner: učitel jako postavička z modelíny ve stylu 3D
animovaného filmu, pozadí: zelená tabule s nápisem 'Do toho!',
světlé barvy, čitelný nápis uprostřed, vhodné pro slide 16:9."
3) Nálepka/štítek pro úkol:
"Samolepka 'Hotovo' s bílým okrajem, kruhový tvar, barvy: žlutá + černá,
plochý ikonický styl, vysoký kontrast pro tisk na samolepky."
\end{verbatim}

Ukázka postupu pro úpravu části obrázku (transformace)
\begin{enumerate}
  \item Vygenerujte základní obrázek v Designeru nebo nahrajte vlastní fotografii.
  \item V Copilot chatu zvolte režim Vytvořit + Transformace, popište přesnou změnu (např. „vyměň pozadí za kreslenou zelenou tabuli a přidej nápis 'Do toho!'“). Pozn.: tato funkce je dostupná u Microsoft 365 Copilot (placená) \cite{kb_ai_101_tipu_pdf_04e4a115_page_52_2}.
  \item Klepněte na Aktualizovat a opakujte doladění, dokud výsledek neodpovídá potřebám hodiny \cite{kb_ai_101_tipu_pdf_04e4a115_page_52_2}.
\end{enumerate}

Praktické nápady použití ve výuce (rychlé příklady)
\begin{itemize}
  \item Vytvoření sady ikon/štítků pro typy aktivit (cvičení, domácí úkol, checkpoint) pomocí Designeru a export do balíku ikon pro slide šablony \cite{kb_ai_101_tipu_pdf_04e4a115_page_51_1}.
  \item Generování ilustračních bannerů pro začátek lekce (motivace, cíle) a jejich iterativní úprava podle zpětné vazby studentů prostřednictvím transformací \cite{kb_ai_101_tipu_pdf_04e4a115_page_52_2}.
  \item Příprava materiálů pro výtvarnou výuku: monogramy nebo varianty v různých malířských technikách (praktické příklady v Designeru) \cite{kb_ai_101_tipu_pdf_04e4a115_page_51_1}.
\end{itemize}

Checklist při exportu a nasazení vizuálů (stručně; general background knowledge)
\begin{itemize}
  \item Zkontrolujte rozlišení a poměr stran pro cílové médium (projekce, tisk, web).
  \item Přidejte stručný alt‑text k obrázkům a poznámku pro lektora, jak obrázek v hodině komentovat.
  \item Uložte originální prompty a metadata (styl, poměr stran), aby bylo možné vizuály reprodukovat nebo upravovat později. 
  \item Pokud používáte materiály studentů v publikaci, ověřte souhlas a zásady ochrany osobních údajů. 
\end{itemize}

Poznámka k licencím a dostupnosti funkcí
\begin{itemize}
  \item Některé pokročilé editační funkce (např. transformace v Copilot chatu) vyžadují placenou licenci Microsoft 365 Copilot; také některé šablony/obsahy v Designeru mohou být omezeny pro typy účtů \cite{kb_ai_101_tipu_pdf_04e4a115_page_52_2,kb_ai_101_tipu_pdf_04e4a115_page_51_1}.
\end{itemize}